\chapter{Evaluation of Approaches}

This appendix contains information about the evaluation of the two approaches. We show analytically the calculated metric values.


%===================================
\section{Fine-Tuned LLM Results}

The following table demonstrates the calculated values of the evaluation metrics for the fine-tuned LLM approach.


\begin{table}[h!]
    \centering
    \caption{Evaluation metrics across different training epochs}
    \label{tab:results_ft}
    \resizebox{\textwidth}{!}{%
    \begin{tabular}{lcccccc}
    \hline
    \textbf{Epochs} & \textbf{BERT Precision} & \textbf{BERT Recall} & \textbf{BERT F1} & \textbf{Match@3} & \textbf{Precision@3} & \textbf{Recall@3} \\
    \hline
    3 Epochs & 0.9229 & 0.9133 & 0.9181 & 0.7862 & 0.3145 & 0.3145 \\
    5 Epochs & 0.9228 & 0.9202 & 0.9215 & 0.7799 & 0.3187 & 0.3187 \\
    8 Epochs & 0.9188 & 0.9179 & 0.9183 & 0.6981 & 0.2851 & 0.2851 \\
    \hline
    \end{tabular}%
    }
\end{table}



%===================================
\section{RAG Pipeline Results}

The following table demonstrates the calculated values of the evaluation metrics for the RAG pipeline approach.

\begin{table}[h!]
    \centering
    \caption{Evaluation metrics across different approaches}
    \label{tab:results_rag}
    \resizebox{\textwidth}{!}{%
    \begin{tabular}{lcccccc}
    \hline
    \textbf{Attempts} & \textbf{BERT Precision} & \textbf{BERT Recall} & \textbf{BERT F1} & \textbf{Match@3} & \textbf{Precision@3} & \textbf{Recall@3} \\
    \hline
    10 Justifications & 0.8727 & 0.8381 & 0.8550 & 0.8302 & 0.3669 & 0.3669 \\
    15 Justifications & 0.8648 & 0.8314 & 0.8477 & 0.7862 & 0.3548 & 0.3522 \\
    20 Justifications & 0.8608 & 0.8290 & 0.8446 & 0.7107 & 0.2778 & 0.2893 \\
    \hline
    \end{tabular}%
    }
\end{table}


%===================================
\section{Approaches Comparison}

The following table demonstrates the calculated values of the evaluation metrics for the fine-tuned LLM and RAG pipeline approaches. The best attempt of each system is selected. 


%===================================
\subsection{Semantic \& Recommendation Metrics}

Below we compare the metrics that measure semantic similarity and recommendation accuracy.

\begin{table}[h!]
    \centering
    \caption{Semantic \& Recommendation Metrics across both Approaches}
    \label{tab:results_ft_rag}
    \resizebox{\textwidth}{!}{%
    \begin{tabular}{lcccccc}
    \hline
    \textbf{Approach} & \textbf{BERT Precision} & \textbf{BERT Recall} & \textbf{BERT F1} & \textbf{Match@3} & \textbf{Precision@3} & \textbf{Recall@3} \\
    \hline
    Fine-Tuned LLM & 0.9228 & 0.9202 & 0.9215 & 0.7799 & 0.3187 & 0.3187 \\
    RAG Pipeline & 0.8727 & 0.8381 & 0.8550 & 0.8302 & 0.3669 & 0.3669 \\
    \hline
    \end{tabular}%
    }
\end{table}

%===================================
\subsection{Efficiency Metrics}

Below, we compare the metrics that measure the runtime duration of each approach.

\begin{table}[h!]
    \centering
    \caption{Efficiency Metrics across both Approaches}
    \label{tab:results_ft_rag_time}
    \begin{tabular}{lcc}
    \hline
    \textbf{Approach} & \textbf{Average Time per Query (s)} & \textbf{Total Runtime (mm:ss)} \\
    \hline
    Fine-Tuned LLM & 16.00 & 12:31 \\
    RAG Pipeline   & 37.56 & 42:24 \\
    \hline
    \end{tabular}
\end{table}


%===================================
\subsection{Strengths and Limitations}

Below, we provide a table that analytically describes the advantages and disadvantages of the two frameworks. 


\begin{longtable}{>{\bfseries}p{0.20\textwidth} Y Y}
    \caption{Fine-tuned LLM vs.\ RAG with Re-Ranking} \\
    \label{tab:comparison_ft_rag}
    \toprule
    Aspect & Fine-tuned LLM & RAG with Re-Ranking Pipeline \\
    \midrule
    \endfirsthead
    
    \toprule
    Aspect & Fine-tuned LLM & RAG with Re-Ranking Pipeline \\
    \midrule
    \endhead
    
    \midrule
    \multicolumn{3}{r}{\emph{Continued on next page}} \\
    \midrule
    \endfoot
    
    \bottomrule
    \endlastfoot
    
    Core Idea
    & Train an LLM (Mistral-7B-Instruct with QLoRA) to directly generate top-3 business location recommendations with justifications.
    & Use embeddings to retrieve relevant neighborhoods and re-rank based on similarity to the user query; LLM (Llama 3) only formats the final output. \\
    
    Strengths
    & \begin{itemize}[leftmargin=*, nosep]
        \item Fast inference after training (if GPU is utilized)
        \item Learns domain-specific patterns
        \item Can generalize to unseen but similar prompts
      \end{itemize}
    & \begin{itemize}[leftmargin=*, nosep]
        \item Dynamic: no need to retrain when data updates
        \item Reduces hallucinations (data fetched at runtime)
        \item Easier to personalize results to user preferences
      \end{itemize} \\
    
    Limitations
    & \begin{itemize}[leftmargin=*, nosep]
        \item Static: retraining required for new NACE classes or updated proxy truths
        \item Risk of hallucination in justifications
        \item High initial computational cost
        \item One-to-one conversation does not fully utilize the GPU
      \end{itemize}
    & \begin{itemize}[leftmargin=*, nosep]
        \item Slower inference (retrieval {+} embedding overhead)
        \item Dependent on embedding quality
        \item Requires additional infrastructure (vector DB / storage)
      \end{itemize} \\
    
    Explainability
    & Built-in through feature \& neighborhood justification (SHAP {+} percentile templates).
    & Explainability derived from neighborhood descriptions (feature statistics \& contextual text). \\
    
    Adaptability
    & Less flexible -- retraining needed when data or categories change.
    & Highly flexible -- can add new businesses, neighborhoods, or user constraints without retraining. \\
    
    Computational Resources
    & Requires HPC (A100 GPU) for fine-tuning; optimized via QLoRA (0.18\% trainable params).
    & Moderate GPU/CPU needed for embeddings and inference; feasible to run locally (quantized models). \\
    
    Personalization
    & Limited -- produces 3 static recommendations per NACE class.
    & Strong -- query-specific re-ranking allows user-tailored recommendations. \\
    
    Best Use Case
    & When fast inference on fixed, stable data with concurrent queries or batched workloads is needed, and compute resources are available.
    & When frequent updates, user-specific preferences, and reduced hallucinations are priorities. \\
    
\end{longtable}


